\Article{365}
It is enough for us here to have shown the possibility of these solutions. For brevity, we pass over many devices that can ease these operations. But we cannot omit some very important consequences.

What precedes proves that all the auxiliary equations for solving the equation $x^\nu=1$, if converted into congruences, have possible roots when the period
\SourceLeaf{244}
\[
x+x^p+x^{pp}+\dotstwo+x^{p^{m-1}}
\]
has not yet been split. Let us stay with the case where $\nu$ is prime; then $m$ is a divisor of $\nu-1$. Here, therefore, the auxiliary congruences have real roots if the number of periods found through them is a divisor of $\frac{\nu-1}m$. Thus, if $\frac{\nu-1}m$ is even, that is, if $m$ is a divisor of $\frac{\nu-1}2$, or equivalently if $p^{\frac{\nu-1}2}\equiv1\;(\mathrm{mod.}\,\nu)$, or equivalently if $p$ is a quadratic residue of the prime $\nu$, the quadratic equation that divides the roots into two periods will have real roots modulo $p$. But in Chapter \textsc{vi} we showed that, on putting $\nu=4n\pm1$, this equation is always $xx+x\mp n=0$. Hence we have the notable

\Needspace{12\baselineskip}\textsc{Theorem.} \emph{If the prime $p$ is a quadratic residue of the prime $4n\pm1$, the congruence}
\[
xx+x\mp n\equiv0(\mathrm{mod.}\,p)
\]
\emph{will have real roots, and therefore so will the congruence}
\[
4xx+4x\mp4n\equiv0\quad\text{or}\quad(2x+1)^2\mp\nu\equiv0
\]
\emph{that is, $\pm\nu$ will be a quadratic residue of $p$.}

\Article{366}
This, then, is the third complete proof of the fundamental theorem of Chapter \textsc{iv}, all the more worthy of attention because the principles from which it is drawn are entirely different from those we used for the earlier proofs.\ .But from this same source, by the opposite route, let us derive a fourth. Namely, let $\nu$ be a prime of the form $4n\pm1$, let $p$ be any other prime, and let $\pm\nu$ be a quadratic residue of the prime $p$; we shall prove that $p$ is a quadratic residue of $\nu$.

Let $p^m$ be the least power of $p$ that is $\equiv1(\mathrm{mod.}\,\nu)$. The elementary divisors of the function $\frac{x^\nu-1}{x-1}$ modulo $p$ will have degree $m$, so their total number will be $=\frac{\nu-1}m$. Now, since $\pm\nu Rp$, the congruence
\[
xx+x\mp n\equiv0(\mathrm{mod.}\,p)
\]
will be solvable; let its roots be $A,A'$. Distribute the functions $x,xx,\dotstwo x^{\nu-1}$ into two classes, denoted by $\xi,\xi'$. We have
\SourceLeaf{245}
\[
\begin{aligned}
\xi+\xi'&\equiv A+A'+(1+x+xx+\dotstwo+x^{\nu-1})\\
\xi\xi'&\equiv AA'+\lambda(1+x+xx+\dotstwo+x^{\nu-1})
\end{aligned}
\]
so that
\[
(z-\xi)(z-\xi')-(z-A)(z-A')
\]
will be divisible by every elementary divisor of the function $\frac{x^\nu-1}{x-1}$. Each of these elementary divisors will therefore divide either $\xi-A$ and $\xi'-A'$, or $\xi-A'$ and $\xi'-A$. It follows (since $A$ is not $\equiv A'$) that, if in place of $x$ we substitute $x^p$, then $\xi$ and $\xi'$ do not change. For if $\xi$ passed into $\xi'$, and conversely, then $\xi-A$ and $\xi-A'$ would be divisible by the same prime function. Q.\ E.\ A. Finally, it follows that $\frac{\nu-1}2$ is divisible by $m$, or equivalently that $p^{\frac{\nu-1}2}-1$ is divisible by $\nu$. Thus $p$ is a quadratic residue of $\nu$. Q.\ E.\ D.

It is easy to derive all cases of the fundamental theorem from the two theorems.

\Article{367}
Although we have restricted ourselves here to the case where $\nu$ is prime, analogous theorems can also be determined without much difficulty when $\nu$ is composite; brevity does not now allow us to set this out more fully.

It is evident that similar observations can also be made about a greater number of periods. Thus, if $\frac{\nu-1}m$ is divisible by $3$, that is, if $p$ is a cubic residue of the prime $\nu$, the equation that distributes the roots of the equation $x^\nu=1$ into three periods, and whose determination a priori we taught in Chapter \textsc{vi}, will be solvable modulo $p$, and conversely. For example, the congruence $x^3+xx-2x-1\equiv0$ can be solved modulo any prime of the form $7n\pm1$, but cannot if it has another form.

It would not be difficult for us to enrich this chapter with many other observations, if the limits within which we must keep did not prevent it. For those who wish to go further, these principles can at least point the way.
